\begin{center}
    {\LARGE \textbf{2002无机化学}} \\[2ex]
    {\large 时量180分钟 \quad 满分100分}
\end{center}

\vspace{0.5cm}
\noindent\textbf{注意事项:}
\begin{enumerate}[label=\arabic*., parsep=0pt, itemsep=0pt]
    \item 答题前,考生先将自己的姓名、学号、考场号、座位号填写在答题卡上,将条形码准确粘贴在考生信息条形码粘贴区。
    \item 考试结束后,将本试卷与答题纸一并收回。
    \item 回答各个问题时,将答案写在答题纸上,写在本试卷上无效。
\end{enumerate}

\vspace{0.5cm}
\noindent\textbf{一、按要求回答下列问题，每题10分，共10小题，共100分。}

\begin{enumerate}[label=\arabic*., start=1, itemsep=1.5ex]
    \item 完成并配平下列化学反应方程式:
    \begin{enumerate}[label=(\arabic*), itemsep=1ex]
        \item 向\ce{FeSO4}溶液中加入酸性\ce{K2Cr2O7}溶液。
        \item 用盐酸处理\ce{Co(OH)3}
        \item 硝酸银受热分解
        \item 向\ce{HgCl2}溶液中滴加\ce{SnCl2}溶液
        \item \ce{PbS}中加入过量\ce{H2O2}
    \end{enumerate}

    \item \ce{M^{3+}}离子的\ce{3d}轨道中电子半充满。试推断
    \begin{enumerate}
        \item \ce{M}原子的核外电子排布式
        \item \ce{M}元素所在的族和周期
        \item \ce{M}的原子序数和元素符号
    \end{enumerate}

    \item 用分子轨道理论说明氦的双原子分子不能存在的原因，并讨论\ce{He^{+}_{2}}存在的可能性

    \item 标准氢电极\ce{2H^{+} + 2e^{-} -> H_{2}}，$\phi^{\ominus}_{A} = 0\ \text{V}$，若将电极中的标准酸换成$\text{pH}=4$的酸，求$\phi$。如果温度、电极反应\ce{2H_{2}O + 2e^{-} -> H_{2} + 2OH^{-}}，$\phi^{\ominus}_{B} = -0.80\ \text{V}$，试求该温度下水的离子积常数$K_\text{w}$的值。

    \item 已知\ce{[CuCl4]^{2-}}为正方形结构，试用配位化合物价键理论解释之，并对这种解释中不令人满意之处加以说明。

    \item 下列物质能用来鉴别\ce{Na2SO4}溶液和\ce{Na2SO3}溶液的是？请说明实验现象以证实你的选择。
    \begin{tasks}(4)
        \task 盐酸
        \task \ce{NaOH}溶液
        \task \ce{AgCl}
        \task 碘水
    \end{tasks}

    \item \ce{Ag^{+}},\ce{Hg^{2+}},\ce{Zn^{2+}},\ce{Al^{3+}},\ce{Cr^{3+}}离子共存，试设计一个不生成硫化物的分离方案。

    \item 以海水为主要原料，使用\ce{Cl2},\ce{Na2CO3}等试剂制备单质溴。试叙述其过程，并写出反应的化学方程式。

    \item 将无色硝酸盐\ce{A}放入水中，生成白色沉淀物\ce{B}和澄清的溶液\ce{C}。向其中通入\ce{H2S}生成黑色沉淀\ce{D}，\ce{D}不溶于\ce{NaOH}溶液，但可溶于盐酸。向\ce{C}中滴加\ce{NaOH}溶液，有白色沉淀\ce{E}生成，\ce{E}不溶于过量的\ce{NaOH}溶液。试给出\ce{A}、\ce{B}、\ce{C}、\ce{D}、\ce{E}的化学式。

    \item 氨元素的氧化物有多种。试写出其化学式，并分别描述其物理性质及化学性质。

\end{enumerate}